\section{Curvas Características del Transistor}

Los transistores son dispositivos no lineales. Su comportamiento cambia drásticamente dependiendo de cuánta corriente o voltaje se les inyecte. A diferencia de resistores o componentes lineales, no se los puede definirlo con una sola ecuación simple; más bien se requiere un mapa de su comportamiento. Ese mapa son las curvas características.

A continuación verá que un transistor tiene cuatro regiones importantes. Cada región representa un modo de operación completamente distinto. Dependiendo de lo que se esté diseñando o reparando, se fuerza el componente a trabajar en una región específica.

Para que esta sección no sea extremadamente tediosa, Carl ayudará a Luis a comprender el significado de cada región.

\subsection{Curva característica de entrada}
Observando el circuito de la figura \ref{fig_circuito_em_comm} ¿Qué forma cree que tiene la gráfica de la corriente de base \(I_B\) en función de la tensión base-emisor \(V_{BE}\) (\(I_B(V_{BE})\))? Preguntó Carl. Si recuerda, esta unión se comporta exactamente como un diodo convencional. Entonces la respuesta es inmediata: su gráfica es idéntica a la de un diodo polarizado en directa.
\begin{figure}[!ht]
  \centering
  \begin{circuitikz}[scale=0.8,transform shape]
    \draw (0,0) node[npn](Q){};
    \draw (Q.C) to[short] ++(0,0.5) to[R,l={\(R_C\)}] ++(2.5,0) to[battery,l={\(V_{CC}\)}] ++(0,-3) coordinate(GND) node[ground]{};
    \draw (Q.B) to[R,l_={\(R_B\)}] ++(-3,0) coordinate(tmp) to[battery,l_={\(V_{BB}\)}] (GND -| tmp) node[ground]{};
    \draw (Q.E) -- (GND -| Q.E) node[ground]{};
  \end{circuitikz}
  \caption{Topología de circuito con configuración EC.}
  \label{fig_circuito_em_comm}
\end{figure}

Por lo tanto, podemos aplicar las mismas simplificaciones estudiadas anteriormente en el capítulo \ref{chpt_introduccion}. Al analizar la malla de entrada, la corriente se define por la ley de Ohm:
\[
  I_B = \frac{V_{BB} - V_{BE}}{R_B}
\]
Para el análisis práctico de circuitos, la segunda aproximación del diodo es el equilibrio perfecto entre simplicidad y precisión. Basta con asumir que, una vez encendido el transistor, la barrera de potencial \(V_{BE}\) se mantiene casi constante en aproximadamente \qty{0.7}{\volt}.

\subsection{Curva de colector y regiones de operación}

La verdadera naturaleza del transistor se revela al observar su salida. ``¿A qué llamamos salida?'', interrumpe Luis.

``Llamamos salida a la corriente de colector en función de la tensión colector-emisor'', respondió Carl, y anotó en la pizarra
\[
  I_C(V_{CE}) \quad \text{(Función de salida)}
\]

Para poder estudiar esta función, supongamos que fijamos la corriente de base \(I_B\) en un valor pequeño y constante (por ejemplo, \qty{10}{\micro\ampere}). A continuación, comenzamos a variar la tensión de la fuente del colector y medimos tanto la corriente de colector \(I_C\) como la tensión en los terminales \(V_{CE}\). 

Al trazar estas mediciones, la gráfica resultante no es una línea recta como dictaría una resistencia simple, sino que revela distintas zonas de comportamiento dictadas por el diseño interno del componente.

La figura \ref{fig_curvas_de_colector} representa las curvas para distintas corrientes de base (\qty{10}{\micro\ampere}, \qty{20}{\micro\ampere}, etcétera).

\begin{figure}[!ht]
  \centering
  \begin{tikzpicture}[scale=0.8,>=stealth]
  \draw[->,gray] (-.4,0) -- (8.4,0) node[right,font=\scriptsize]{\(V_{CE}\)};
  \draw[->,gray] (0,-.4) -- (0,8.4) node[above,font=\scriptsize]{\(I_C\) [\unit{\milli\ampere}]};

  \draw[thick,teal!40] (0,0) -- ++(20:.05) to[out=20,in=180] ++(.35,.1) -- 
    node[midway,above,font=\tiny,color=black]{\(I_B=\qty{0}{\ampere}\) (Corriente de base)} ++(7.5,0) to[out=0,in=260] ++(.2,.5) 
    -- ++(.2,1.8);
  \draw[thick,teal!40] (0,0) -- ++(70:.5) to[out=70,in=180] ++(.35,.3) -- 
    node[midway,above,font=\tiny,color=black]{\qty{10}{\micro\ampere}} ++(7.1,0) to[out=0,in=260] ++(.2,.5) 
    -- ++(.2,2);
  \foreach \i [evaluate=\i as \corriente using int(20*\i)] in {1,1.5,2,2.5,3,3.5} {
    \draw[thick,teal!40] (0,0) -- ++(70:\i) to[out=70,in=180] ++(.35,.4) -- 
      node[midway,above,font=\tiny,color=black]{\qty{\corriente}{\micro\ampere}} ++(7.6-\i,0) to[out=0,in=260] ++(.2,.5) 
      -- ++(.2,{\i+1/\i});
  }
  \foreach \i in {1,2,3,4,5,6,7,9,10,11,12,13,14,15}{
    \draw[gray,thin] (.1,0.52*\i+.2) -- (-.1,0.52*\i+.2) node[left,font=\tiny]{\qty{\i}{\milli\ampere}};
  }
  \draw[gray,thin] (1,0) -- (1,-0.1) node[below,font=\tiny]{\qty{1}{\volt}};
  \draw[gray,thin] (7.6,0) -- (7.6,-0.1) node[below,font=\tiny]{\qty{40}{\volt}};
  \draw[gray!50!blue,->] (-2.3,5) node[above,text width=1.5cm,align=center,font=\scriptsize,draw]{Región de saturación} to[out=-90,in=120] (70:4);
  \draw[gray!50!red,->] (10,5) node[above,text width=1.5cm,align=center,font=\scriptsize,draw]{Región de disrupción} to[out=-90,in=60] (8,4);
  \draw[black!50!gray!50!green,->] (3.5,6.5) node[above,text width=1cm,align=center,font=\scriptsize,draw]{Región activa} -- ++(0,-2);
\end{tikzpicture}

  \caption{Conjunto de curvas de colector.}
  \label{fig_curvas_de_colector}
\end{figure}

\paragraph{La región activa}
Si observa la gráfica a medida que \(V_{CE}\) supera un valor mínimo, verá que la curva se vuelve casi completamente horizontal. Es decir, \(I_C\) se mantiene constante aunque sigamos aumentando la tensión entre el colector y el emisor. 

``¿Por qué ocurre esto?'', preguntó Carl. A lo que Luis, recordando la explicación física que antes dedujo, respondió: ``El colector se limita a atrapar los electrones que el emisor logra inyectar a través de la base. La cantidad de electrones disponibles depende única y exclusivamente de la corriente de base \(I_B\). Por más que aumentemos el voltaje \(V_{CE}\) para intentar succionar más corriente, el colector no puede absorber electrones que no existen''. Carl complementó esta explicación, añadiendo que en esta zona plana, conocida como región activa, el transistor actúa como un amplificador perfecto donde la corriente de salida es simplemente un múltiplo de la entrada (\(I_C = \beta I_B\)).

\paragraph{La región de saturación}
¿Qué sucede en la parte inicial de la gráfica, cuando \(V_{CE}\) se acerca a \qty{0}{\volt}? En este punto, la tensión del colector es tan débil que su campo eléctrico pierde la fuerza necesaria para capturar a los electrones que cruzan la base. Al no haber un ``tobogán'' lo suficientemente inclinado para barrerlos, la corriente de colector cae abruptamente hacia cero. A esta pendiente casi vertical se la conoce como región de saturación.

\paragraph{La región de disrupción}
Como regla general, el transistor jamás debe operar en esta región, ya que el resultado es su destrucción casi instantánea. Esta zona crítica se alcanza cuando la tensión $V_{CE}$ se eleva a niveles no soportados por el dispositivo. Al hacerlo, el campo eléctrico en la unión Base-Colector se vuelve tan intenso que acelera a los electrones a velocidades catastróficas. Estos chocan contra la estructura cristalina del silicio y arrancan violentamente a otros electrones de sus enlaces, generando un efecto multiplicador de avalancha. El resultado es una corriente masiva y descontrolada que atraviesa el dispositivo, ignorando por completo cualquier orden dictada por la base, y que termina destruyendo el componente de forma irreversible.

\paragraph{La región de corte}
Si realizamos el experimento forzando \(I_B = 0\) (por ejemplo, desconectando la base), esperaríamos que la corriente de colector fuera exactamente cero. En la gráfica, esto se representa como una curva microscópica pegada al eje horizontal inferior. Esta ínfima circulación se debe a corrientes de fuga térmicas propias de los semiconductores, pero en la inmensa mayoría de las aplicaciones prácticas es tan pequeña que el transistor se considera un interruptor abierto perfecto.

\subsection{Familia de curvas y límites de potencia}
Si repetimos el procedimiento midiendo \(I_C\) para distintos valores fijos de \(I_B\) (\qty{20}{\micro\ampere}, \qty{30}{\micro\ampere}, etc.), obtendremos una familia de curvas. Cada peldaño horizontal superior representa una corriente de colector proporcionalmente mayor. Este conjunto de gráficas es la huella dactilar del transistor y define su ganancia de corriente continua (\(\beta\)).

Sin embargo, el funcionamiento en estas regiones tiene límites físicos ineludibles. Aplicando la ley de las tensiones de Kirchhoff en la malla de salida, la caída de potencial en el dispositivo es:
\[
V_{CE} = V_{CC} - I_C R_C
\]
Y la potencia que el transistor debe disipar en forma de calor es el producto de su tensión y su corriente:
\[
P_D = V_{CE} I_C
\]
Si la combinación de voltaje y corriente genera una potencia (\(P_D\)) que supera la potencia máxima especificada por el fabricante, la temperatura interna excederá rápidamente los \qty{150}{\degreeCelsius} y el cristal de silicio se destruirá por exceso térmico. 

De manera similar, si continuamos aumentando \(V_{CE}\) desproporcionadamente en la gráfica, llegaremos a un muro infranqueable: la \textit{región de disrupción} (o ruptura). Por ello, el diseño de cualquier circuito debe asegurar que el transistor opere siempre dentro de las zonas seguras (activa, corte o saturación), lejos de estos límites destructivos.


